\section{Lean formalization: completed components and a path to completion}
\label{sec:formalization}

\subsection{The two proof layers}

There are two distinct uses of Lean in this project.  At the semantic layer,
\texttt{Nat}, \texttt{Int}, \texttt{Rat}, and lists provide convenient
implementations of the finite mathematics.  Proving a theorem at this
layer checks its mathematical statement in Lean's ambient foundations.
At the object layer, \texttt{PR}, \texttt{PTm}, \texttt{Fm}, and
\texttt{Prf PRAAx} describe the terms, formulas, and derivations of
\ipra.  A completed object proof must use the permitted induction
constructors and appropriate encoded formulas.  Moving between these
layers requires construction of a derivation, not a change of notation.

The standalone \texttt{ObjectKernel.lean} in this archive contains the
declarations and proofs of the pinned upstream \texttt{Syntax.lean},
\texttt{Deduction.lean}, and \texttt{PR.lean}.  Their imports are replaced
by \texttt{Std}, and the natural-number notation is supplied explicitly.
The declaration blocks were compared byte for byte with the originals.
This makes the foundational example reproducible without downloading the
full Mathlib dependency tree.  Upstream licensing and exact hashes are
recorded alongside it.  It is not intended to be imported simultaneously
with the original modules, which define the same names.

\subsection{Actual compiled proof objects}

The central object interface has the following shape; the complete scoped
definition is in \texttt{ObjectBridge.lean}:
\begin{quote}\small\ttfamily
pre : Fm PTm n\\
program : PR n\\
post : Fm PTm (n + 1)\\
certified : Prf PRAAx [pre] (post with output replaced by program)
\end{quote}
The \texttt{certified} field is an actual object derivation.  It is not a
Boolean check of a semantic proposition.  The compiled existential
introduction lemma turns this into an object derivation of
\(\mathrm{pre}\to\exists y\,\mathrm{post}\).

\texttt{ObjectExamples.lean} defines addition using the literal PR syntax:
its primitive recursion starts with a projection, and its step composes
successor with the previous-value projection.  It proves the zero and
successor equations inside \texttt{PRAAx}, and then constructs an actual
object induction proof of
\[
                         \forall x\ (x+0=x).
\]
The induction formula is the equation itself.  This small example is
useful because it exercises the exact scoping, substitution, induction,
and proof-object interfaces that larger applications will use.

The semantic modules also contain useful general lemmas:
\begin{center}\small
\begin{longtable}{@{}P{.27\linewidth}P{.36\linewidth}P{.28\linewidth}@{}}
\toprule
Module & Content & Status\\
\midrule\endhead
\texttt{FiniteContracts} & Bounded iteration invariants; least bounded
search; finite-prefix extensionality; producer composition. & General
semantic proofs compiled.\\
\texttt{Metastability} & Finite rise alternative; bounded change;
constant windows for bounded monotone natural sequences. & General
semantic proofs compiled.\\
\texttt{PowerGap} & Natural gap inequalities; signed integer power gap;
small-perturbation exclusion. & General semantic proofs compiled.\\
\texttt{ObjectKernel}, \texttt{ObjectBridge}, \texttt{ObjectExamples}
& Deep arithmetic syntax and deduction; contract interface; genuine
object induction example. & Compiled object-level infrastructure and
example.\\
\texttt{SeriesTargets} & Truncated rational arrays, norms, fixed-point
iteration, certificate checks, coefficient-name targets. & Definitions,
definitional lemmas, and closed examples compiled; main targets unproved.\\
\texttt{PowerGapTargets} & Polynomial list operations, Boolean certificate
check, height functions, main targets. & Definitions compiled; polynomial
theorem targets unproved.\\
\texttt{AnalyticTargets} & Regular names and order, finite arithmetic
tables, soft endpoint selection, and common analytic targets. & Definitions
and target statements compiled; main targets unproved.\\
\texttt{ChartProducerTargets} & Gaussian-rational data, explicit chart
constants, and simple-root certificate producer targets. & Definitions
compiled; producer proof targets unproved.\\
\texttt{SelectionTargets} & Coherent finite witness selection and
bounded-cardinality empty-window constructions. & Definitions compiled;
main targets unproved.\\
\texttt{RepositoryBridge} & Adapter to the original
\texttt{pra\_iff\_native} interface. & Source outline inspected; not built
against the full original dependency closure.\\
\bottomrule
\end{longtable}
\end{center}

The compiler used is Lean 4.32.0, commit
\texttt{8c9756b28d64dab099da31a4c09229a9e6a2ef35}.
The recorded audits distinguish a compiled proposition definition from a
proved proposition.  They also distinguish Lean's metatheoretic axioms
from the rules occurring in an embedded arithmetic proof.  Some compiled
proofs use \texttt{propext} and \texttt{Quot.sound}; the audited new proof
declarations do not use \texttt{sorryAx} or \texttt{Classical.choice}.
This fact does not by itself establish an \ipra derivation; the type and
constructors of the object proof establish that claim.

\subsection{The common missing compiler work}

The broad framework requires a reusable library of syntactic derivations.
The following implementation sequence makes each missing step concrete.
\begin{enumerate}
\item Define named PR symbols for integer coding, pairing, rational
cross-comparison, list lookup, finite folds, and powers.  Prove their
equations and elementary algebraic laws using \texttt{PRA.qfInd}.
Record the domain conditions as Boolean equations.
\item Define the syntax of bounded formulas over a \emph{fixed} finite
signature, with bounds stored as terms.  Construct the Booleanization
function and an object derivation of the equivalence for every constructor.
The quantified cases must use the guarded-prefix proof in
Theorem~\ref{logic:booleanization}, not induction on the formula with an
unbounded quantifier.
\item Define finite annotated programs: assignment, composition,
Boolean branch, and bounded loop.  An annotation includes its bounded
invariant and already constructed base/step derivations.  Compile the
program to its PR symbol and to its postcondition derivation by
Theorem~\ref{logic:compilation}.  A checker for the resulting raw proof
objects should verify syntax and local inference rules without globally
normalizing arbitrary PR terms.
\item Build fixed certified lemmas for bounded search, finite maximum,
minimum with deterministic tie breaking, weighted finite sums, and
finite convolution.  Each becomes a reusable proof constructor rather
than an assumed semantic property.
\item Introduce the name formulas as ordinary first-order abbreviations.
Every represented theorem should expose the finite precisions it uses,
then universally generalize the requested accuracy only after completing
the finite derivation.  The rational separation lemma eliminates any
remaining limiting phrase in a rational conclusion.
\end{enumerate}
The semantic proofs already supplied are convenient guides for these
proof constructors, but they do not automatically produce them.

\subsection{Completion route for each main mathematical result}

For \textbf{finite precision propagation}, use a list of arithmetic gates
whose argument indices precede the gate.  Define exact and rounded
evaluation by a bounded fold.  Store the maximum previous error as an
auxiliary rational value, and use
\(e_j\le(2B+3)^j\delta\), together with the operand bounds, as the
Booleanized invariant.  Prove the product difference identity by ring
algebra and its inequality by signed integer cases.  This yields the
gate-error target and, after reification, its object derivation.

For \textbf{integration}, define grid indices as vectors bounded by
\(2^{\mu(k+3)}\), flattened to a finite list.  Prove finite refinement
counting and equality of the total weights to one.  A weighted-error
fold gives the exact rational comparison of two grid sums.  Regularity
and modulus formulas are used by explicit universal elimination;
they are not the induction formula.  The continuous extension at named
points uses the separately specified clipping and precision schedule.

For \textbf{metastability}, the natural change-budget theorem is already
proved.  The rational theorem adds the PR floor operation and its
interval property.  The represented theorem can be proved independently:
construct the candidate endpoints and first minimum of rational
differences; prove the finite minimum inequality and the telescoping
identity; insert the two endpoint-error bounds.  The exact rational
margin \(2^{-k}/(2K)\) supplies the output apartness witness.  The
empty-window theorem adds bounded search and the invariant of distinct
stored witnesses.  It uses the supplied cardinality premise only after
constructing the potentially forbidden finite list.

For \textbf{coherent witness selection}, define finite lists \(D_s\),
the PR residual evaluator, the first minimum, and
\(\tau(k)=\max(k,\sigma(k))\).  Prove the two-error estimate for that
minimum.  Instantiate the isolation contract at the two selected points,
and then the continuity contract at the resulting fast name.  No
choice or compactness theorem is needed.

For \textbf{the implicit branch}, prove the six principal finite targets
in \texttt{SeriesTargets.lean} in dependency order: convolution norm,
finite power difference, contraction, invariant ball, geometric
telescoping, and coefficient estimate.  Fix degree and horizon before
every induction.  Add the explicit clipping maps for scalar coefficients
and weighted arrays; their error bound is twice the sample error.
Use the gate lemma to compare two finite polynomial computations.
The coefficient-name formula, its finite weighted bound, coefficientwise
fixed-point equation, and explicit tail then follow.  Finally prove the
finite residual bound for evaluating the series; this avoids assuming
an infinitary substitution theorem.  Reify the same finite targets before
transporting the resulting object proofs.  The supplied rational-array
module is the real scalar prototype; Gaussian pairs use the explicitly
defined pair arithmetic and the norm $|u|+|v|$.  The finite convolution
and contraction proofs use only the displayed norm inequalities, so
the same proof structure applies to those pairs.  The certificate
producer then needs the finite derivative estimates $M_1,M_2$, the
explicit functions $U,h,T,p,W,t$, and the sampled upper-bound estimates
of Corollary~\ref{ana:producer}.  These are additional named targets,
not consequences of defining a certificate structure.

For \textbf{power-gap certificates}, the signed integer gap is already
proved semantically.  Formalize polynomial evaluation of a coefficient
identity, the leading-term and integer-root bounds, and the finite loss
estimate.  Those lemmas prove the outside-box contradiction and hence
the two height bounds.  Bounded enumeration is a generic verified filter.
The power-completion constructor needs its triangular coefficient
invariant and a positive denominator product.  The zero-remainder branch
and the parity analysis for the explicit family are finite algebraic
arguments.  For denominator transport, formalize the three homogenized
coefficient lists and their coefficient-norm estimates.  All of these
algorithms and target predicates are specified in the archive.

\subsection{Integration into ProveIt}

A natural integration location is a reusable library such as
\path{Logic/PrimitiveRecursiveArithmetic/FiniteContracts/}, with the
application reports pointing to it.  This is a proposed organization,
not a repository change made by this task.  In an actual integration,
replace the standalone \texttt{ObjectKernel} import by
\texttt{ConciseIPRA.PR}; do not load both copies.  Build the native
adapter against \texttt{ConciseIPRA.NativeBiInterpretation}.  Its last
step should have the form
\[
 h:\operatorname{Prf}\,\mathrm{PRAAx}\,\Gamma\,\varphi
 \quad\longmapsto\quad
 (\texttt{pra\_iff\_native}).\texttt{mp}\ h.
\]
The arithmetic proof \(h\) is the substantive input, and the equivalence
does not manufacture it.

Once the Booleanization, rational arithmetic, and bounded-loop constructors
are in place, the same work benefits a large fraction of the collection:
finite coefficient identities, exact geometric certificates, contour
grids, truncated series, arithmetic sums, height searches, and finite
topological tables all need these operations.  This is the principal
practical advantage of one common framework over separate ad hoc
arithmetization sections.
